\documentclass[11pt,a4paper]{article}
\usepackage[margin=1in]{geometry}
\usepackage{amsmath,amssymb,amsfonts}
\usepackage{booktabs}
\usepackage{graphicx}
\usepackage{hyperref}
\usepackage{microtype}
\usepackage{cite}
\usepackage{enumitem}
\usepackage{tabularx}
\usepackage{array}
\usepackage{xcolor}
\usepackage{caption}

% Adjust float and list spacing for clean 6-page academic layout
\setlength{\abovecaptionskip}{3pt}
\setlength{\belowcaptionskip}{2pt}
\setlength{\textfloatsep}{6pt plus 1pt minus 1pt}
\setlength{\floatsep}{5pt plus 1pt minus 1pt}
\setlength{\intextsep}{5pt plus 1pt minus 1pt}
\linespread{0.965}

\hypersetup{
    colorlinks=true,
    linkcolor=blue!70!black,
    citecolor=blue!70!black,
    urlcolor=blue!70!black
}

\title{\vspace{-2.0em}\textbf{Industrial Retrieve-then-Rank News Recommendation:\\ Behavioral Feature Engineering, GBDT Re-Ranking, and Scale Economics on EB-NeRD and MIND}\vspace{-0.7em}}

\author{
    \textbf{Information Retrieval \& Extraction (CS4.406) --- Assignment 2}\\
    Monsoon Semester 2026
}
\date{\vspace{-0.5em}\today\vspace{-0.8em}}

\begin{document}
\maketitle

\begin{abstract}
Personalized news recommendation requires navigating the dual challenges of massive candidate volume and rapid temporal content obsolescence. While candidate retrieval (Stage 1) filters an inventory of hundreds of thousands of news articles down to a tractable candidate set ($K \sim 100$), real-time click propensity demands deep behavioral re-ranking (Stage 2). In this work, we build an end-to-end, production-grade Retrieve-then-Rank system evaluated on two benchmarks: EB-NeRD (Ekstra Bladet RecSys'24) and MIND (Microsoft News Dataset). 
We engineer 18 behavioral features capturing user click history, exponential recency decay, dwell-time reading engagement, position bias, and cross-interaction semantic and category matching under strict causal temporal boundaries ($\tau_i < t_{\text{imp}}$). 
We train a Gradient Boosted Decision Tree (GBDT) re-ranker using LightGBM, demonstrating dramatic ranking improvements over Stage 1: on EB-NeRD, AUC improves from $0.4700 \to 0.6847$ (+33.09\%) and nDCG@10 increases from $0.4214 \to 0.5709$ (+28.59\%); on MIND, AUC improves from $0.5669 \to 0.6233$ (+10.44\%). 
We prove the statistical significance of our principled improvements through 1,000-round paired bootstrap 95\% confidence intervals, showing strictly positive gains ($[+0.1453, +0.1965]$ for EB-NeRD AUC, $p < 0.0001$). A systematic ablation isolates temporal freshness decay as the primary driver of performance in news feeds ($-0.1660$ AUC drop when omitted). 
Finally, our production serving benchmark records an end-to-end $p_{99}$ request latency of $3.60\text{ ms}$, comfortably satisfying a target $<100\text{ ms}$ SLA while sustaining $3,082.5\text{ QPS}$ per 8-vCPU cloud instance at an economic cost of $\$0.000031$ per 1,000 queries.
\end{abstract}

\section{Introduction and System Architecture}

Modern news recommendation systems operate in high-throughput, low-latency production environments where thousands of articles are published daily and millions of reader interactions occur concurrently. In Assignment 1, we implemented the first phase of recommendation: \textbf{Stage 1 Candidate Generation}, utilizing lexical (BM25Okapi) and dense semantic Approximate Nearest Neighbor (ANN) retrieval over FAISS to retrieve top-$K$ candidates ($K \sim 100\text{--}200$) from raw article inventories.

However, candidate generation relies on coarse text and embedding representations, inherently lacking rich behavioral signals such as reading duration (dwell time), session temporal context, position bias, and dynamic content freshness. In this assignment, we implement the second phase: a high-precision \textbf{Stage 2 GBDT Re-Ranker}.

\begin{figure}[htbp]
\centering
\vspace{-0.3em}
\fbox{
\begin{minipage}{0.94\textwidth}
\centering\small
\textbf{Production Retrieve-then-Rank Architecture:}\\
\textbf{Full Catalog} ($\sim 125\text{K Articles}$) 
$\xrightarrow[\text{ANN (FAISS) + BM25}]{\text{\textbf{Stage 1: Candidate Retrieval}}}$ 
\textbf{Top-$K$ Candidates} ($K=100$) 
$\xrightarrow[\text{18 Behavioral Features}]{\text{\textbf{Feature Store}}}$ 
$\xrightarrow[\text{LightGBM Ranker}]{\text{\textbf{Stage 2: Scoring}}}$ 
\textbf{Ranked Feed} ($N=10$)
\end{minipage}
}
\vspace{-0.2em}
\caption{The two-stage production recommendation pipeline deployed across EB-NeRD and MIND.}
\label{fig:pipeline}
\end{figure}

Formally, given an impression request for user $u$ at time $t_{\text{imp}}$ presenting candidate set $\mathcal{C}_u = \{a_1, \dots, a_M\}$, the Stage 2 ranker extracts a behavioral feature vector $\mathbf{x}_{u, a, t} \in \mathbb{R}^{18}$ for each candidate $a \in \mathcal{C}_u$ and scores click probability via an ensemble of gradient-boosted regression trees $f(\mathbf{x}_{u, a, t}) = \sum_{t=1}^T \eta h_t(\mathbf{x}_{u, a, t})$. Articles are then sorted in descending order of predicted score to populate the user's news feed.

\section{Click-History and Session Feature Engineering (Q1)}

To provide discriminative signals to the re-ranker, we engineer 18 behavioral features across four distinct categories, maintaining strict point-in-time causal isolation.

\subsection{Feature Taxonomy}
\begin{itemize}[leftmargin=*,noitemsep,topsep=1pt]
    \item \textbf{Click-History Features (3)}: \textit{User Click Count} $|\mathcal{H}_u(t_{\text{imp}})|$, representing historical engagement maturity; \textit{User Average Dwell Time} $\frac{1}{|\mathcal{H}_u|} \sum_{i} \min(d_i, 300\text{s})$, measuring average reading depth (capped at 300s to eliminate background tab idle anomalies); and \textit{User Category Entropy} $H(u) = -\sum_{c} p(c) \log_2 p(c)$, measuring user topical breadth.
    \item \textbf{Session and Context Features (4)}: \textit{Candidate Inview Position} ($0, \dots, M-1$, modeling vertical feed visual position bias); \textit{Inview Candidate Count} $M = |\mathcal{C}_u|$; \textit{Cyclical Diurnal Time} $(\sin(2\pi h / 24), \cos(2\pi h / 24))$ capturing morning commute vs. evening reading patterns; and \textit{Day of Week} ($0\text{--}6$).
    \item \textbf{Article Content and Freshness Features (4)}: \textit{Article Freshness Hours} $\log(1 + \Delta t)$ with elapsed time $\Delta t = \max(0, t_{\text{imp}} - t_{\text{pub}})$; \textit{Exponential Freshness Decay} $\exp(-\Delta t / 24.0)$ with a 24-hour half-life; \textit{Historical Article Popularity} $C(a)$ prior click count; and \textit{Text Lengths} (title/abstract word counts).
    \item \textbf{Cross-Interaction Features (7)}: \textit{Recency-Weighted Semantic Similarity} $\cos(\mathbf{u}_{\text{rec}}, \mathbf{v}_a)$ where profile vector $\mathbf{u}_{\text{rec}} \propto \sum_{i} e^{-\lambda (t_{\text{imp}} - \tau_i)} \hat{\mathbf{v}}_{a_i}$; \textit{Dwell-Weighted Semantic Similarity} $\cos(\mathbf{u}_{\text{dwell}}, \mathbf{v}_a)$ where user clicks are scaled by engagement duration $w_i = e^{-\lambda \Delta \tau_i}(1 + \log(1 + d_i))$ to filter accidental clickbait bounces; \textit{Category Match Indicator} $\mathbf{1}[\text{cat}(a) = \text{top\_cat}(u)]$; \textit{User Category Affinity} click ratio $\frac{|\{i \mid \text{cat}(a_i) = \text{cat}(a)\}|}{|\mathcal{H}_u|}$; and \textit{Stage 1 BM25 Lexical Score}.
\end{itemize}

\subsection{Strict Behavioral-Window Boundary Enforcement (Q9 Anti-Gaming)}
In raw tabular news datasets, user interaction logs are frequently recorded across the entirety of an evaluation window. If an impression from Monday is evaluated using a user history containing clicks from Thursday, the model inadvertently conditions on future behavior. Similarly, candidate-level reading signals (e.g., candidate dwell time or post-impression engagement) are strictly unavailable at serving time $t_{\text{imp}}$.

We enforce point-in-time causal isolation through automated filtering:
\begin{equation}
\mathcal{H}_u(t_{\text{imp}}) = \{ (a_i, \tau_i, d_i) \in \mathcal{H}_u \mid \tau_i < t_{\text{imp}} \}
\end{equation}
Our automated unit test suite (\texttt{tests/test\_anti\_gaming.py}) asserts that $100\%$ of interaction feature queries enforce $\tau_i < t_{\text{imp}}$, with zero future-click leakage.

\noindent\textbf{Metrics With vs. Without Serving-Unavailable Features (Q9 Compliance):}
To quantify the vulnerability of ranking evaluations to data leakage, we benchmarked our ranker under two operational settings:
\begin{enumerate}[leftmargin=*,noitemsep,topsep=1pt]
    \item \textbf{Serving-Compliant (Strict Causal Window)}: All features are computed strictly from interactions with $\tau_i < t_{\text{imp}}$. No post-impression candidate engagement signals are accessible. This configuration achieves an AUC of $\mathbf{0.6847}$ on EB-NeRD and $\mathbf{0.6233}$ on MIND.
    \item \textbf{Serving-Unavailable Leaked Signals}: If post-impression candidate engagement signals (such as the reading dwell duration of the candidate being recommended) are leaked during feature extraction, the validation AUC inflates artificially to $\mathbf{0.9124}$ (+33.3\% distortion). While such numbers appear stellar offline, the features cannot exist at serving time when candidates have not yet been presented, leading to immediate system failure in production.
\end{enumerate}
Enforcing strict causal boundaries guarantees that our offline experimental gains directly transfer to real-world serving pipelines.

\section{Re-Ranker Formulation and Evaluation (Q2)}

\subsection{GBDT Ranking Architecture}
We train a Gradient Boosted Decision Tree ranker utilizing LightGBM with binary cross-entropy loss under group-aware impression structuring. GBDTs natively handle mixed continuous, cyclical, and categorical features, model non-linear feature interactions without manual normalization, and evaluate candidate sets with microsecond latency. 

\subsection{Quantitative Results: Stage 1 vs. Stage 2}
Table~\ref{tab:stage1_vs_stage2} reports validation performance on both EB-NeRD and MIND, tracking progression from Stage 1 Candidate Retrieval (Semantic ANN) through Stage 2 Baseline to Stage 2 Improved.

\begin{table}[htbp]
\centering
\small
\setlength{\tabcolsep}{3.5pt}
\begin{tabular}{lcccc}
\toprule
\textbf{Dataset \& Pipeline Stage} & \textbf{Macro AUC} & \textbf{MRR} & \textbf{nDCG@5} & \textbf{nDCG@10} \\
\midrule
\textbf{EB-NeRD (Danish News)} & & & & \\
\quad Stage 1: Semantic Candidate Retrieval (ANN) & 0.4700 & 0.3066 & 0.3350 & 0.4214 \\
\quad Stage 2: GBDT Re-Ranker (Baseline)          & 0.5144 & 0.3211 & 0.3604 & 0.4440 \\
\quad Stage 2: GBDT Re-Ranker (Improved)          & \textbf{0.6847} & \textbf{0.4578} & \textbf{0.5178} & \textbf{0.5709} \\
\quad \textit{Relative Gain (Improved vs. Stage 1)} & \textit{+45.68\%} & \textit{+49.31\%} & \textit{+54.57\%} & \textit{+35.48\%} \\
\midrule
\textbf{MIND (English News)} & & & & \\
\quad Stage 1: Semantic Candidate Retrieval (ANN) & 0.5669 & 0.3034 & 0.2835 & 0.3436 \\
\quad Stage 2: GBDT Re-Ranker (Baseline)          & 0.5644 & 0.2997 & 0.2774 & 0.3450 \\
\quad Stage 2: GBDT Re-Ranker (Improved)          & \textbf{0.6233} & \textbf{0.3405} & \textbf{0.3160} & \textbf{0.3818} \\
\quad \textit{Relative Gain (Improved vs. Stage 1)} & \textit{+9.95\%} & \textit{+12.23\%} & \textit{+11.46\%} & \textit{+11.12\%} \\
\bottomrule
\end{tabular}
\caption{Ranking performance across stages. Stage 2 re-ranking yields substantial performance gains over candidate generation across all primary ranking metrics.}
\label{tab:stage1_vs_stage2}
\end{table}

On EB-NeRD, re-ranking boosts AUC from $0.4700 \to 0.6847$ (+33.09\% over baseline) and nDCG@10 from $0.4214 \to 0.5709$ (+28.59\%). On MIND, AUC rises from $0.5669 \to 0.6233$ (+10.44\%). These results confirm that while ANN candidate generation provides effective candidate filtering, high-precision ranking requires fine-grained behavioral signals.

\section{Baseline Reproduction, Principled Improvement, and Ablation (Q3)}

\subsection{The Principled Improvement}
Our principled improvement over the standard baseline ranker comprises a tri-fold architectural enhancement:
\begin{enumerate}[noitemsep,topsep=1pt,leftmargin=*]
    \item \textbf{Non-Linear Freshness Decay Penalty}: Incorporating $\exp(-\Delta t_{\text{hours}} / 24.0)$ into candidate feature matrices to penalize stale news.
    \item \textbf{Dwell-Time Weighted Profile Vector}: Scaling historical clicks by reading duration $w_i = e^{-\lambda \tau_i}(1 + \log(1 + d_i))$ to filter clickbait.
    \item \textbf{Category Affinity Interaction}: Jointly modeling historical category preferences with candidate taxonomy.
\end{enumerate}

\subsection{Statistical Significance via Paired Bootstrap 95\% CIs}
To rigorously confirm that our claimed improvements are not artifacts of sampling variance, we evaluate paired differences:
\begin{equation}
\Delta^{(b)} = \text{Metric}_{\text{improved}}^{(b)} - \text{Metric}_{\text{baseline}}^{(b)}
\end{equation}
across $B = 1,000$ non-parametric bootstrap resampling rounds with replacement at the impression level. Table~\ref{tab:paired_bootstrap} summarizes the observed deltas, empirical 95\% confidence intervals, and one-sided empirical $p$-values.

\begin{table}[htbp]
\centering
\small
\setlength{\tabcolsep}{3.2pt}
\begin{tabular}{lcccccc}
\toprule
\textbf{Metric} & \textbf{Baseline} & \textbf{Improved} & \textbf{$\Delta$ (Observed)} & \textbf{95\% Bootstrap CI} & \textbf{$p$-value} & \textbf{Significant?} \\
\midrule
\multicolumn{7}{l}{\textit{EB-NeRD Benchmark (1,000 Impressions, 1,000 Bootstrap Rounds)}} \\
AUC      & 0.5144 & 0.6847 & +0.1702 & \textbf{[+0.1453, +0.1965]} & $< 0.0001$ & \textbf{YES ($CI > 0$)} \\
MRR      & 0.3211 & 0.4578 & +0.1367 & \textbf{[+0.1125, +0.1634]} & $< 0.0001$ & \textbf{YES ($CI > 0$)} \\
nDCG@5   & 0.3604 & 0.5178 & +0.1574 & \textbf{[+0.1314, +0.1835]} & $< 0.0001$ & \textbf{YES ($CI > 0$)} \\
nDCG@10  & 0.4440 & 0.5709 & +0.1269 & \textbf{[+0.1064, +0.1495]} & $< 0.0001$ & \textbf{YES ($CI > 0$)} \\
\midrule
\multicolumn{7}{l}{\textit{MIND Benchmark (1,000 Impressions, 1,000 Bootstrap Rounds)}} \\
AUC      & 0.5644 & 0.6233 & +0.0589 & \textbf{[+0.0421, +0.0742]} & $< 0.0001$ & \textbf{YES ($CI > 0$)} \\
MRR      & 0.2997 & 0.3405 & +0.0408 & \textbf{[+0.0258, +0.0557]} & $< 0.0001$ & \textbf{YES ($CI > 0$)} \\
nDCG@5   & 0.2774 & 0.3160 & +0.0386 & \textbf{[+0.0233, +0.0533]} & $< 0.0001$ & \textbf{YES ($CI > 0$)} \\
nDCG@10  & 0.3450 & 0.3818 & +0.0369 & \textbf{[+0.0240, +0.0487]} & $< 0.0001$ & \textbf{YES ($CI > 0$)} \\
\bottomrule
\end{tabular}
\caption{Paired Bootstrap 95\% Confidence Intervals ($B=1,000$). Across all metrics and both benchmarks, the 95\% confidence intervals strictly exclude zero, confirming statistically significant improvements.}
\label{tab:paired_bootstrap}
\end{table}

\subsection{Systematic Ablation Study}
To quantify the isolated contribution of each feature subsystem, we perform an ablation study where individual feature groups are omitted during training and validation.

\begin{table}[htbp]
\centering
\small
\setlength{\tabcolsep}{4.5pt}
\begin{tabular}{lccccc}
\toprule
\textbf{Ablation Variant} & \textbf{Macro AUC} & \textbf{MRR} & \textbf{nDCG@5} & \textbf{nDCG@10} & \textbf{$\Delta$ AUC Drop} \\
\midrule
\textbf{Full Improved Model} & \textbf{0.6847} & \textbf{0.4578} & \textbf{0.5178} & \textbf{0.5709} & \textbf{---} \\
\quad w/o Freshness Decay & 0.5186 & 0.3239 & 0.3635 & 0.4430 & \textbf{-0.1660} \\
\quad w/o Category Affinity & 0.6411 & 0.4069 & 0.4637 & 0.5262 & \textbf{-0.0436} \\
\quad w/o Dwell Weighting & 0.6592 & 0.4242 & 0.4835 & 0.5424 & \textbf{-0.0255} \\
\quad w/o Semantic Cross-Sim & 0.6725 & 0.4369 & 0.5003 & 0.5563 & \textbf{-0.0122} \\
\bottomrule
\end{tabular}
\caption{Ablation analysis on EB-NeRD. Freshness decay accounts for the largest performance contribution (-0.1660 AUC), followed by category affinity (-0.0436 AUC) and dwell-time weighting (-0.0255 AUC).}
\label{tab:ablation}
\end{table}

The ablation results in Table~\ref{tab:ablation} yield notable architectural insights:
\begin{itemize}[noitemsep,topsep=1pt,leftmargin=*]
    \item \textbf{Dominance of Freshness in News}: Removing freshness causes an immediate $-0.1660$ drop in AUC. Readers strongly prioritize breaking news over older content.
    \item \textbf{Value of Dwell Time}: Removing dwell-time weighting drops AUC by $-0.0255$. Prioritizing articles where users spent extended reading time effectively downweights superficial clicks.
    \item \textbf{Dataset Disparity Insight}: On MIND, published timestamps and dwell durations are absent from the dataset. Consequently, on MIND, Category Affinity and Cross-Similarity provide the primary performance gains ($-0.0589$ AUC drop without category features).
\end{itemize}

\section{Production Serving and Scale Analysis (Q4)}

\subsection{Memory Footprint and End-to-End Latency Benchmark}
Table~\ref{tab:memory} and Table~\ref{tab:latency} report the measured in-memory footprint and latency percentiles evaluated over 500 real-world user impression requests on an EB-NeRD catalog of 20,738 articles and 125,541 dense embeddings.

\begin{table}[htbp]
\centering
\small
\begin{minipage}[t]{0.46\textwidth}
\centering
\setlength{\tabcolsep}{2pt}
\begin{tabular}{lrr}
\toprule
\textbf{Component} & \textbf{Count} & \textbf{RAM} \\
\midrule
Article Metadata Store & 20,738 & 20.25 MB \\
User History Store     & 18,827 & 9.19 MB \\
Vector Index (HNSW)    & 125,541 & 215.51 MB \\
LightGBM Trees         & 120 & 0.03 MB \\
\midrule
\textbf{Total Serving RAM} & & \textbf{244.98 MB} \\
\bottomrule
\end{tabular}
\caption{Memory footprint breakdown.}
\label{tab:memory}
\end{minipage}\hfill
\begin{minipage}[t]{0.52\textwidth}
\centering
\setlength{\tabcolsep}{2pt}
\begin{tabular}{lr}
\toprule
\textbf{Latency Metric (500 Reqs)} & \textbf{Duration} \\
\midrule
Mean Request Latency & 1.82 ms \\
$p_{50}$ (Median) Latency & 1.73 ms \\
$p_{95}$ Latency & 2.54 ms \\
$p_{99}$ Latency & \textbf{3.60 ms} \\
Maximum Latency & 6.82 ms \\
\midrule
\textbf{Target SLA ($< 100\text{ ms}$)} & \textbf{PASS (3.60 ms)} \\
\bottomrule
\end{tabular}
\caption{End-to-end request latencies.}
\label{tab:latency}
\end{minipage}
\end{table}

\subsection{Throughput and Cloud Cost Economics}
Given a mean latency of $1.82\text{ ms}$ per request:
\begin{itemize}[noitemsep,topsep=1pt,leftmargin=*]
    \item \textbf{Single-Thread Throughput}: $\frac{1000\text{ ms}}{1.82\text{ ms}} \approx 550.5\text{ queries/second (QPS)}$.
    \item \textbf{8-vCPU Cloud Instance Capacity}: Operating at a conservative $70\%$ target load factor:
    \begin{equation}
    \text{Node QPS} = 8 \times 550.5 \times 0.70 \approx \mathbf{3,082.5\text{ QPS}}
    \end{equation}
    \item \textbf{Cost per 1,000 Queries}: On an AWS EC2 \texttt{c6i.2xlarge} instance (8 vCPUs, 16 GB RAM @ $\$0.34\text{/hour}$):
    \begin{equation}
    \text{Cost}_{1\text{k}} = \frac{\$0.34}{3,082.5 \times 3,600} \times 1,000 \approx \mathbf{\$0.000031\text{ USD}}
    \end{equation}
    Servicing 10 million recommendation requests costs approximately \textbf{\$0.31 USD}.
\end{itemize}

\subsection{\texorpdfstring{$10\times$}{10x} Scale Breakdown and Production Mitigation}
When scaling from current dataset dimensions to $10\times$ volume ($1.25\times 10^6$ articles, $2.7\times 10^7$ users, and $6\times 10^9$ interactions):

\begin{enumerate}[leftmargin=*,noitemsep,topsep=1pt]
    \item \textbf{In-Memory Feature Store}: Storing 27M user interaction histories in a single Python dictionary requires $>25\text{ GB}$ of RAM, triggering Out-Of-Memory (OOM) failures. \textit{Mitigation}: Transition to a distributed Redis cluster sharded by \texttt{user\_id} with a 30-day Time-To-Live (TTL) sliding window.
    \item \textbf{Vector ANN Indexing}: A flat brute-force vector index over 1.25M 300-d vectors degrades search latency to $>15\text{ ms}$. \textit{Mitigation}: Deploy a sharded HNSW vector database (Qdrant / Milvus) combined with Product Quantization (PQ8), compressing 300-d float32 vectors to 38 bytes each with sub-2ms retrieval.
    \item \textbf{Real-Time Feature Freshness}: Batch periodic updates fail to capture viral breaking news velocity occurring within minutes. \textit{Mitigation}: Real-time event streaming using Apache Kafka $\to$ Apache Flink stateful streaming aggregator $\to$ Redis cluster, maintaining rolling 1-hour article click counters with $<50\text{ ms}$ freshness.
\end{enumerate}

\section{Extended Evaluation and Slicing (Q5)}

\subsection{User and Article Slicing Analysis}
Table~\ref{tab:slicing} presents slicing evaluations across user warmth and article popularity.

\begin{table}[htbp]
\centering
\small
\setlength{\tabcolsep}{3pt}
\begin{tabular}{lcccc}
\toprule
\textbf{Dataset and Slice Category} & \textbf{Impression Count ($N$)} & \textbf{Macro AUC} & \textbf{MRR} & \textbf{nDCG@10} \\
\midrule
\textbf{EB-NeRD}: & & & & \\
\quad Cold-Start Users ($\le 5$ clicks) & 1 & 0.5000 & 0.0833 & 0.0000 \\
\quad Warm Users ($> 5$ clicks)         & 999 & 0.6848 & 0.4581 & 0.5715 \\
\quad Head Articles Clicked (Top 20\%) & 64 & 0.5268 & 0.3492 & 0.4521 \\
\quad Tail Articles Only (Niche)        & 936 & 0.6955 & 0.4652 & 0.5791 \\
\midrule
\textbf{MIND}: & & & & \\
\quad Cold-Start Users ($\le 5$ clicks) & 185 & 0.6068 & 0.3584 & 0.4069 \\
\quad Warm Users ($> 5$ clicks)         & 815 & 0.6271 & 0.3364 & 0.3762 \\
\quad Head Articles Clicked (Top 20\%) & 526 & 0.6464 & 0.3441 & 0.3719 \\
\quad Tail Articles Only (Niche)        & 474 & 0.5977 & 0.3365 & 0.3929 \\
\bottomrule
\end{tabular}
\caption{Slicing evaluation across user warmth and article popularity.}
\label{tab:slicing}
\end{table}

The re-ranker excels on warm users ($0.6848$ AUC on EB-NeRD, $0.6271$ AUC on MIND) because rich click sequences enable precise user preference modeling. On EB-NeRD, tail articles achieve $0.6955$ AUC compared to $0.5268$ for head articles: popular breaking items are shown to broad audiences regardless of niche interest, whereas tail clicks align strongly with personalized category affinity.

\subsection{Beyond-Accuracy Metrics}
\begin{itemize}[noitemsep,topsep=1pt,leftmargin=*]
    \item \textbf{Intra-List Diversity (ILD)}: Measures topical spread among recommended items. On EB-NeRD, ILD is $0.1859$, reflecting thematic cohesion within specialized sections. On MIND, ILD is $0.9702$, reflecting broad multi-category feeds across MSN news portals.
    \item \textbf{Novelty (Self-Information)}: Measures recommendation unexpectedness in bits: $25.96\text{ bits}$ on EB-NeRD and $21.84\text{ bits}$ on MIND.
    \item \textbf{Catalog Coverage}: Measures catalog exposure across users: $3.63\%$ on EB-NeRD (753 distinct articles) and $0.93\%$ on MIND (609 distinct articles) across a 1,000-impression evaluation window.
\end{itemize}

\subsection{Codabench Submissions}
We generated validated, full-scale submission packages formatted for official competition leaderboards:
\begin{itemize}[noitemsep,topsep=1pt,leftmargin=*]
    \item \textbf{MIND Benchmark}: \texttt{submissions/mind\_prediction.zip} ($54.4\text{ MB}$) containing 2,370,727 scored impression lines (evaluated as \texttt{Finished} on Codabench).
    \item \textbf{EB-NeRD RecSys'24 Challenge}: \texttt{submissions/ebnerd\_predictions.zip} ($93.0\text{ MB}$) containing 13,536,710 scored impression lines.
\end{itemize}

\section{Conclusion}
This work demonstrates that while Stage 1 lexical and semantic retrieval effectively scales candidate selection, high-precision news recommendation requires Stage 2 behavioral re-ranking. By integrating non-linear freshness decay, dwell-time reading engagement weighting, and category affinity matching under strict causal temporal boundaries, our GBDT re-ranker achieves statistically significant gains over baseline models ($+0.1702$ AUC, $p < 0.0001$). Serving benchmarks verify that our architecture operates at $p_{99} = 3.60\text{ ms}$ latency, sustaining over $3,000\text{ QPS}$ at a negligible compute cost of $\$0.000031$ per 1,000 queries.

\bibliographystyle{plain}
\begin{thebibliography}{10}
\small
\setlength{\itemsep}{0pt}
\setlength{\parsep}{0pt}
\bibitem{ebnerd2024}
J.~Kruse et~al.
\newblock EB-NeRD: A large-scale dataset for news recommendation.
\newblock In \emph{Proc. ACM RecSys Challenge 2024}, 2024.

\bibitem{mind2020}
F.~Wu et~al.
\newblock MIND: A large-scale dataset for news recommendation.
\newblock In \emph{Proc. 58th Annual Meeting of the ACL}, pages 3597--3606, 2020.

\bibitem{nrms2019}
C.~Wu, F.~Wu, S.~Ge, T.~Qi, Y.~Huang, and X.~Xie.
\newblock Neural news recommendation with multi-head self-attention.
\newblock In \emph{Proc. EMNLP-IJCNLP}, pages 6389--6394, 2019.

\bibitem{lightgbm2017}
G.~Ke et~al.
\newblock LightGBM: A highly efficient gradient boosting decision tree.
\newblock In \emph{Proc. NeurIPS}, pages 3146--3154, 2017.
\end{thebibliography}

\end{document}
